\chapter{Introdução}
\section{Contextualização e Problematização}
A dregradação da qualidade do ar atmosférico e de ambientes 
internos representar um dos maiores desafios da saúde pública me
escala mundia, estados diretamente associada ao surgimento e ao 
agravamento de enfermidades respiratórias e cardiovasculares.
Segundo estimativas da Organizaçãoi Mundial da Saúde (OMS) a poluição do ar causa
cerca de 4,2 milhões de mortes prematuras por ano. No Brasil, calcula - se
que a degradação da qualidade do ar atmosférico cause aproximadamente 20 mil óbitos
anuais, acrescidos de mais de 10,7 mil mortes por ano decorrentes da poluição do ar
em ambientes internos \cite{cunha2024}. Esse panorama é severamente intensificado
na região amazônica devido ao impacto direto e recorrente das queimadas florestais,
que elevam expressivamente a emissõ de material particulado e gases nocivos na
atmosféra resultando na poluição \cite{gomes2024}.


Apesar dos avanços recentes na expansão da Internet das Coisas (IoT), a maioria
das soluções atuais de sensoriamento de baixo custo opera sob a abordagem puramente reativa 
e passiva \cite{cunha2024}. Nesses sistemas convencionais, o nó sensor limita - se a coltar
medições brutas e transmiti - las à nuvem para exibição em dashboards estátivos, sem capacidade
de antecipar picos de contaminação por poluentes como o Monoxido de Carbono (CO),
Ozonio (O\textsubscript{3}), Dióxido de Carbono(CO\textsubscript{2}) ou Amonia NH\textsubscript{3}), e 
sem executar ações de mitigação imediata. \cite{cunha2024}

Do ponto de vista da arquitetura de sistemas embarcados, a maioria dos nós sensores IoT
é desenvolvido sobre a arquitetura de laço infinito. Nesse modelo
sem escalonamentos preempivo, a CPU do microcontrolador interrompe a rotina de amostragem
física dos sensores durante a montagem e o envio dos pacotes via LoRaWan, introduzindo
jitter, atrasos no tempos de transmissão (alcançando médias atpes 8,05 segundos) e uma taxa
de perda de pacotes de 15,95\%. Alem disso, a dependencia exclusiva de servidores
de nuvem para processar algoritmos de inteligencia artificial gera vulnerabilidade a quedas 
de rede e latencia elevadas, impedindo garantias deterministicas de tempo real em cenários
de riscos a saúde \cite{alves2025freertos}

\section{Justificativa da Pesquisa}